
\section{Conclusión}

Al terminar esta práctica nos dimos cuenta de que no necesitas ser hacker para encontrar bastante información sobre alguien. Con buscar su nombre en internet y revisar algunas redes sociales puedes encontrar muchas cosas, como sus gustos, viajes, estudios y algunas fechas importantes \cite{cwsi_osint}. En nuestro caso, encontramos perfiles de LinkedIn, GitHub, Facebook e Instagram que nos ayudaron a conocer más sobre los dos agentes.

Lo que más nos llamó la atención fue ver cómo cosas que parecen normales, como subir una foto, hablar de una serie o aparecer en un lugar, pueden dar más información de la que pensamos cuando se juntan con otros datos. También vimos que encontrar las ubicaciones de las fotos no siempre es tan fácil, porque las redes sociales suelen borrar sus metadatos y \texttt{exiftool} no nos sirvió para obtener las coordenadas \cite{exiftool}. Aun así, pudimos identificar algunos lugares fijándonos en los edificios y otros detalles que aparecían en las imágenes. También notamos que encontramos mucha más información de Chipotles que de Miri, lo que nos hizo ver que todo depende de lo que cada persona comparte públicamente.

Sobre el caso de Shosholanda, creemos que \textbf{no hay evidencia suficiente} para decir que Miri o Chipotles son culpables o inocentes. Aunque encontramos información sobre sus actividades personales, sus estudios y sus gustos, no encontramos pruebas que demuestren que filtraron contraseñas de la empresa. Esto tampoco significa que no haya pasado, simplemente que con la información que encontramos no podemos asegurarlo. Para saber qué ocurrió realmente se necesitarían pruebas más directas, como los registros internos de la empresa. Toda la investigación se hizo con fines escolares y la información recopilada se utiliza únicamente para este reporte.

Finalmente, esta práctica nos hizo pensar más en lo que publicamos en internet. A veces subimos fotos o compartimos cosas sin pensar que alguien puede juntar esa información y conocer más de nosotros de lo que imaginamos. Por eso, es importante tener cuidado con lo que publicamos, revisar la privacidad de nuestras cuentas y pensar un poco antes de compartir algo, porque aunque después lo borremos, no siempre significa que desaparezca por completo. Esto también nos ayuda a entender cómo se puede aprovechar esa información en ataques de ingeniería social \cite{crowdstrike_social}.
